\section{Tabelle}
\begin{table}[H]
\centering
\caption{Portfolio representations generated by different kernels}
\label{tab:kernel_policy_classes}
\small
\renewcommand{\arraystretch}{1.25}

\begin{tabular}{
p{0.21\textwidth}
p{0.34\textwidth}
p{0.35\textwidth}
}
\toprule
Representation & Kernel & Portfolio policy\\
\midrule

Common exposure
&
$K(z,z')=1$
&
$W(z)$ is constant
\\[4pt]

Linear
&
$K(z,z')=1+z^\prime z'$
&
$W(z)$ is linear in characteristics
\\[4pt]

Finite features
&
$K(z,z')=\Phi(z)^\prime\Phi(z')$
&
$W(z)$ lies in the span of the chosen features
\\[4pt]

Polynomial
&
$K(z,z')=(1+z^\prime z')^q$
&
$W(z)$ includes polynomial terms and interactions
\\[6pt]

Gaussian
&
$\displaystyle
K(z,z')
=
\exp\!\left(
-\frac{\|z-z'\|^2}{2\ell^2}
\right)$
&
Flexible nonlinear $W(z)$
\\[8pt]

Mat\'ern
&
$\displaystyle
\begin{aligned}
K_s(z,z')
&=
c_{s,D}\,
\rho_s(z,z')^{\,s-D/2}
\\[-1pt]
&\quad\times
\mathcal K_{s-D/2}
\!\bigl(\rho_s(z,z')\bigr)
\end{aligned}$
&
Flexible nonlinear $W(z)$ with smoothness controlled by $s$
\\[10pt]

Neural tangent
&
$\displaystyle
\begin{aligned}
K_{\mathrm{NTK}}(z,z')
&=
\frac{\|z\|\,\|z'\|}{2\pi}
\Bigl[\sin\theta
\\[-1pt]
&\qquad
+2(\pi-\theta)\cos\theta
\Bigr]
\end{aligned}$
&
Infinite-width one-hidden-layer ReLU policy
\\

\bottomrule
\end{tabular}

\begin{minipage}{0.97\textwidth}
\footnotesize
\emph{Notes:}
Here $\ell>0$ denotes the kernel length scale. For the Mat\'ern kernel,
\[
\rho_s(z,z')
=
\sqrt{2s-D}\,
\frac{\|z-z'\|}{\ell},
\qquad
c_{s,D}
=
\frac{2^{1-(s-D/2)}}{\Gamma(s-D/2)},
\qquad
s>\frac{D}{2},
\]
and $\mathcal K_\alpha$ denotes the modified Bessel function of the second kind.
Under standard regularity conditions, the Mat\'ern RKHS is equivalent to a
Sobolev space of smoothness $s$.

For the neural tangent kernel,
\[
\theta
=
\arccos\!\left(
\frac{z^\prime z'}
{\|z\|\,\|z'\|}
\right),
\]
and the displayed expression corresponds to a one-hidden-layer ReLU network
without bias in the infinite-width NTK limit, up to the conventional overall
normalization.

The kernel determines the class of characteristic-to-exposure mappings
available to the portfolio manager. Portfolio returns use the common
cross-sectional normalization in Equation~\eqref{eq:portfolio_return}.
The common-exposure specification is an equal-exposure characteristic rule,
not an unrestricted vector of static stock weights.
\end{minipage}
\end{table}


\Needspace{15\baselineskip}
Table~\ref{tab:special-rates} relates the learning exponent to familiar
spectral benchmarks. These are conditions on the managed-portfolio operator,
not conclusions implied by the kernel name alone.

\begin{table}[H]
\centering
\caption{Benchmark portfolio-learning rates}
\label{tab:special-rates}
\small
\begin{tabular}{lcc}
\toprule
Portfolio class
&
Managed-portfolio spectrum
&
Sharpe learning rate
\\
\midrule
Finite dimensional
&
Fixed finite rank
&
$T^{-1}$
\\[5pt]

Gaussian benchmark
&
Fast decay
&
Nearly $T^{-1}$
\\[5pt]

Sobolev--Mat\'ern benchmark
&
$b=2\widetilde s/D$
&
$T^{-2\widetilde s/(2\widetilde s+D)}$
\\
\bottomrule
\end{tabular}

\par\smallskip
\begin{minipage}{0.97\textwidth}
\footnotesize
\emph{Notes:} Sharpe rates use Condition~\ref{ec:sharpe}.
The finite-rank rate uses a fixed identifiable payoff space and root-$T$
moment estimation. The fast-decay row is a separate i.i.d. or sufficiently
fast-mixing benchmark, with possible logarithmic factors and an admissible
regularization sequence; it does not follow by setting $b=\infty$ in
Theorem~\ref{thm:rate}. The Sobolev--Mat\'ern row assumes the displayed
\emph{managed-payoff} spectral law and $\gamma>b$.
\end{minipage}
\end{table}
\subsection*{Population-optimal portfolio policy in a factor economy}

\noindent\textbf{Proposition.}
Consider the three-factor return model
\[
R_{i,t+1}=\beta(Z_{i,t})^\top F_{t+1}+\varepsilon_{i,t+1},\qquad i=1,\ldots,N,
\]
where $\beta(z)\in\mathbb R^3$ and the characteristics $Z_{i,t}$ are measurable with respect to $\mathcal F_t$. Assume that
\[
\begin{aligned}
\mathbb E[F_{t+1}\mid\mathcal F_t]&=\mu_F,
&\mathbb E[F_{t+1}F_{t+1}^\top\mid\mathcal F_t]&=I_3,\\
\mathbb E[\varepsilon_{t+1}\mid\mathcal F_t]&=0,
&\mathbb E[\varepsilon_{t+1}\varepsilon_{t+1}^\top\mid\mathcal F_t]&=\sigma_\varepsilon^2I_N,\\
\mathbb E[F_{t+1}\varepsilon_{t+1}^\top\mid\mathcal F_t]&=0,
&&\sigma_\varepsilon^2>0.
\end{aligned}
\]
Suppose that the characteristics have a common marginal distribution $P_Z$, that $\beta_1,\beta_2,\beta_3\in\mathcal H_K$, and that the factor loadings satisfy
\[
\frac1N\sum_{i=1}^N\beta(Z_{i,t})\beta(Z_{i,t})^\top=gI_3,\qquad g=\frac{c_\beta}{3}>0,
\]
almost surely at every date. Then the population-optimal policy minimizing
\[
Q(W)=\mathbb E\left[\left(1-\frac1N\sum_{i=1}^NW(Z_{i,t})R_{i,t+1}\right)^2\right]
\]
over $W\in\mathcal H_K$ is given, uniquely up to $P_Z$-almost-sure equality, by
\[
W^\star(z)=\frac{\beta(z)^\top\mu_F}{g+\sigma_\varepsilon^2/N}.
\]

\begin{proof}
For any $W\in\mathcal H_K$, define the portfolio factor exposure
\[
b_t(W)=\frac1N\sum_{i=1}^NW(Z_{i,t})\beta(Z_{i,t})\in\mathbb R^3.
\]
Substituting the factor model into the portfolio payoff yields
\[
R^p_{t+1}(W)=b_t(W)^\top F_{t+1}+\frac1N\sum_{i=1}^NW(Z_{i,t})\varepsilon_{i,t+1}.
\]
Using the conditional moment assumptions and the law of iterated expectations, we obtain
\[
\mathbb E[R^p_{t+1}(W)\mid\mathcal F_t]=b_t(W)^\top\mu_F
\]
and
\[
\mathbb E[(R^p_{t+1}(W))^2\mid\mathcal F_t]
=\|b_t(W)\|^2+\frac{\sigma_\varepsilon^2}{N^2}\sum_{i=1}^NW(Z_{i,t})^2.
\]
Consequently, using the common marginal distribution of the characteristics, the unconditional population loss becomes
\[
Q(W)=1-2\mu_F^\top\mathbb E[b_t(W)]+\mathbb E[\|b_t(W)\|^2]
+\frac{\sigma_\varepsilon^2}{N}\mathbb E[W(Z)^2].
\]

To compute its functional derivative, take an arbitrary direction $V\in\mathcal H_K$ and consider $W_h=W+hV$, with $h\in\mathbb R$. Since $b_t$ is linear, we have $b_t(W+hV)=b_t(W)+hb_t(V)$. Substituting into the loss and expanding gives
\[
\begin{aligned}
Q(W+hV)=Q(W)
&+2h\left[-\mu_F^\top\mathbb E[b_t(V)]
+\mathbb E[b_t(W)^\top b_t(V)]
+\frac{\sigma_\varepsilon^2}{N}\mathbb E[W(Z)V(Z)]\right]\\
&+h^2\left[\mathbb E[\|b_t(V)\|^2]
+\frac{\sigma_\varepsilon^2}{N}\mathbb E[V(Z)^2]\right].
\end{aligned}
\]
Therefore, the Gâteaux derivative of $Q$ at $W$ in the direction $V$ is
\[
\begin{aligned}
DQ(W)[V]
&=\left.\frac{d}{dh}Q(W+hV)\right|_{h=0}\\
&=-2\mu_F^\top\mathbb E[b_t(V)]
+2\mathbb E[b_t(W)^\top b_t(V)]
+\frac{2\sigma_\varepsilon^2}{N}\mathbb E[W(Z)V(Z)].
\end{aligned}
\]
A population minimizer must satisfy $DQ(W)[V]=0$ for every $V\in\mathcal H_K$. Dividing by two and recalling that
\[
b_t(V)=\frac1N\sum_{i=1}^NV(Z_{i,t})\beta(Z_{i,t}),
\]
the functional first-order condition can be written as
\[
0=\mathbb E\left[\frac1N\sum_{i=1}^NV(Z_{i,t})
\left(\beta(Z_{i,t})^\top[b_t(W)-\mu_F]
+\frac{\sigma_\varepsilon^2}{N}W(Z_{i,t})\right)\right].
\]

We solve this equation by imposing the stronger pointwise condition
\[
\beta(Z_{i,t})^\top[b_t(W)-\mu_F]
+\frac{\sigma_\varepsilon^2}{N}W(Z_{i,t})=0
\]
for every $i$ and almost every realization of the characteristics. This condition is sufficient, although not necessarily equivalent to the functional first-order condition, since it makes the expression inside the expectation vanish for every admissible direction $V$.

Isolating $W(Z_{i,t})$, we obtain
\[
W(Z_{i,t})=\frac{N}{\sigma_\varepsilon^2}\beta(Z_{i,t})^\top[\mu_F-b_t(W)].
\]
Substituting this expression into the definition of $b_t(W)$ yields
\[
\begin{aligned}
b_t(W)
&=\frac1N\sum_{i=1}^N\beta(Z_{i,t})W(Z_{i,t})\\
&=\frac1{\sigma_\varepsilon^2}
\left(\sum_{i=1}^N\beta(Z_{i,t})\beta(Z_{i,t})^\top\right)
[\mu_F-b_t(W)].
\end{aligned}
\]
By the cross-sectional balance condition, the matrix inside the parentheses equals $NgI_3$. Hence
\[
b_t(W)=\frac{Ng}{\sigma_\varepsilon^2}[\mu_F-b_t(W)],
\]
which implies
\[
(Ng+\sigma_\varepsilon^2)b_t(W)=Ng\mu_F,
\qquad
b_t(W)=\frac{Ng}{Ng+\sigma_\varepsilon^2}\mu_F.
\]
Substituting back into the expression for $W(Z_{i,t})$, we find
\[
\begin{aligned}
W(Z_{i,t})
&=\frac{N}{\sigma_\varepsilon^2}\beta(Z_{i,t})^\top
\left[\mu_F-\frac{Ng}{Ng+\sigma_\varepsilon^2}\mu_F\right]\\
&=\frac{N}{\sigma_\varepsilon^2}\beta(Z_{i,t})^\top
\frac{\sigma_\varepsilon^2}{Ng+\sigma_\varepsilon^2}\mu_F\\
&=\frac{\beta(Z_{i,t})^\top\mu_F}{g+\sigma_\varepsilon^2/N}.
\end{aligned}
\]
We have thus derived the function
\[
W^\star(z)=\frac{\beta(z)^\top\mu_F}{g+\sigma_\varepsilon^2/N}
=\frac{\beta(z)^\top\mu_F}{c_\beta/3+\sigma_\varepsilon^2/N}.
\]
Its functional form was not imposed in advance. Moreover, since $\beta_1,\beta_2,\beta_3\in\mathcal H_K$, the resulting $W^\star$ belongs to $\mathcal H_K$.

We now verify that this function satisfies the functional first-order condition. Using cross-sectional balance,
\[
\begin{aligned}
b_t(W^\star)
&=\frac1N\sum_{i=1}^N\beta(Z_{i,t})W^\star(Z_{i,t})\\
&=\frac1{g+\sigma_\varepsilon^2/N}
\left(\frac1N\sum_{i=1}^N\beta(Z_{i,t})\beta(Z_{i,t})^\top\right)\mu_F\\
&=\frac{g}{g+\sigma_\varepsilon^2/N}\mu_F.
\end{aligned}
\]
Therefore, for every $i$ and almost every realization of the characteristics,
\[
\begin{aligned}
&\beta(Z_{i,t})^\top[b_t(W^\star)-\mu_F]
+\frac{\sigma_\varepsilon^2}{N}W^\star(Z_{i,t})\\
&\qquad=
-\frac{\sigma_\varepsilon^2/N}{g+\sigma_\varepsilon^2/N}
\beta(Z_{i,t})^\top\mu_F
+\frac{\sigma_\varepsilon^2/N}{g+\sigma_\varepsilon^2/N}
\beta(Z_{i,t})^\top\mu_F=0.
\end{aligned}
\]
It follows that $DQ(W^\star)[V]=0$ for every $V\in\mathcal H_K$.

To establish global optimality, observe that the second functional variation is
\[
D^2Q(W)[V,V]
=2\mathbb E[\|b_t(V)\|^2]
+\frac{2\sigma_\varepsilon^2}{N}\mathbb E[V(Z)^2].
\]
Both terms are nonnegative, and the second is strictly positive whenever $V$ is nonzero in $L^2(P_Z)$. Thus $Q$ is strictly convex up to $P_Z$-almost-sure equality, and the stationary point $W^\star$ is its unique population minimizer in this sense.

More explicitly, for any $W\in\mathcal H_K$, let $\Delta W=W-W^\star$. Since the first functional variation vanishes at $W^\star$, the quadratic expansion gives the exact excess-loss identity
\[
Q(W)-Q(W^\star)
=\mathbb E[\|b_t(\Delta W)\|^2]
+\frac{\sigma_\varepsilon^2}{N}\mathbb E[\Delta W(Z)^2]\geq0.
\]

Finally, the solution also satisfies the conditional pricing restriction. Let $B_t\in\mathbb R^{N\times3}$ contain the factor loadings as its rows and define
\[
w_t^\star=\frac1N(W^\star(Z_{1,t}),\ldots,W^\star(Z_{N,t}))^\top,
\qquad
v_t=\frac1N(V(Z_{1,t}),\ldots,V(Z_{N,t}))^\top.
\]
The pointwise condition established above is equivalent to
\[
(B_tB_t^\top+\sigma_\varepsilon^2I_N)w_t^\star=B_t\mu_F.
\]
Consequently, for every $V\in\mathcal H_K$,
\[
\begin{aligned}
&\mathbb E[(1-R^p_{t+1}(W^\star))R^p_{t+1}(V)\mid\mathcal F_t]\\
&\qquad=v_t^\top B_t\mu_F
-v_t^\top(B_tB_t^\top+\sigma_\varepsilon^2I_N)w_t^\star=0.
\end{aligned}
\]
Hence the population-optimal policy also satisfies Condition E4.
\end{proof}




For the conditional Gaussian specification
\[
F_{t+1}\sim\mathcal N(\mu_F,I_3-\mu_F\mu_F^\top),
\qquad
\varepsilon_{t+1}\sim\mathcal N(0,\sigma_\varepsilon^2I_N),
\]
with conditionally independent factors and idiosyncratic innovations,
and $\|\beta(z)\|^2=c_\beta$.
Then
\[
\mathbb E[R_{i,t+1}^2\mid\mathcal F_t]
=\|\beta(Z_{i,t})\|^2+\sigma_\varepsilon^2
=c_\beta+\sigma_\varepsilon^2.
\]
Proposition~\ref{prop:gaussian_admissible} therefore verifies
Condition~\ref{ec:returns} with
$B_X^2=8\kappa_K(c_\beta+\sigma_\varepsilon^2)$ and $K_4=3$.
For a normalized Mat\'ern kernel, $\kappa_K=1$.



\begin{proposition}[Eigenvalue decay of the managed-return covariance operator]
Consider the factor economy
\[
R_{i,t+1}=\beta(Z_{i,t})^\top F_{t+1}+\varepsilon_{i,t+1},
\qquad i=1,\ldots,N,
\]
where
\[
\mathbb E[F_{t+1}F_{t+1}^\top\mid\mathcal F_t]=I_3,
\qquad
\mathbb E[\varepsilon_{t+1}\varepsilon_{t+1}^\top\mid\mathcal F_t]
=\sigma_\varepsilon^2I_N,
\qquad
\mathbb E[F_{t+1}\varepsilon_{t+1}^\top\mid\mathcal F_t]=0,
\]
with $\sigma_\varepsilon^2>0$. Assume that the characteristics have a common uniform marginal distribution $P_Z$ on $\mathcal Z=[-1,1]^D$ and that the factor loadings satisfy
\[
\frac1N\sum_{i=1}^N
\beta(Z_{i,t})\beta(Z_{i,t})^\top=gI_3,
\qquad g>0,
\]
almost surely at every date. Let $K$ be a Matérn kernel with smoothness parameter $\nu>0$ and fixed positive length scale, and define
\[
X_t=\frac1N\sum_{i=1}^N R_{i,t+1}K(\cdot,Z_{i,t}),
\qquad
\Sigma_{\mathcal H_K,\varepsilon}
=\mathbb E_\varepsilon[X_t\otimes X_t].
\]
Then, for every fixed economy with finite $N$, the positive eigenvalues $\mu_{j,\varepsilon}$ of $\Sigma_{\mathcal H_K,\varepsilon}$ satisfy
\[
c_\mu j^{-b}\leq\mu_{j,\varepsilon}\leq C_\mu j^{-b},
\qquad j\geq1,
\]
where $b=1+2\nu/D$ and $0<c_\mu\leq C_\mu<\infty$. In particular, for $D=6$ and $\nu=3/2$, we obtain $b=3/2$.
\end{proposition}

\begin{proof}
Consider an arbitrary function $f\in\mathcal H_K$. By the definition of the covariance operator and the reproducing property of the kernel,
\[
\begin{aligned}
\langle f,\Sigma_{\mathcal H_K,\varepsilon}f\rangle_{\mathcal H_K}
&=\mathbb E_\varepsilon[\langle f,X_t\rangle_{\mathcal H_K}^2]\\
&=\mathbb E_\varepsilon\left[
\left(\frac1N\sum_{i=1}^N f(Z_{i,t})R_{i,t+1}\right)^2
\right].
\end{aligned}
\]
Substituting the factor model, we obtain
\[
\frac1N\sum_{i=1}^N f(Z_{i,t})R_{i,t+1}
=\left(\frac1N\sum_{i=1}^N f(Z_{i,t})\beta(Z_{i,t})\right)^\top F_{t+1}
+\frac1N\sum_{i=1}^N f(Z_{i,t})\varepsilon_{i,t+1}.
\]
Since the characteristics are $\mathcal F_t$-measurable, the conditional second-moment assumptions imply
\[
\begin{aligned}
&\mathbb E_\varepsilon\left[
\left(\frac1N\sum_{i=1}^N f(Z_{i,t})R_{i,t+1}\right)^2
\middle|\mathcal F_t\right]\\
&\qquad=
\left\|\frac1N\sum_{i=1}^N f(Z_{i,t})\beta(Z_{i,t})\right\|^2
+\frac{\sigma_\varepsilon^2}{N^2}\sum_{i=1}^N f(Z_{i,t})^2.
\end{aligned}
\]
Taking unconditional expectations and using the common marginal distribution of the characteristics gives
\[
\begin{aligned}
\langle f,\Sigma_{\mathcal H_K,\varepsilon}f\rangle_{\mathcal H_K}
&=\mathbb E_\varepsilon\left[
\left\|\frac1N\sum_{i=1}^N f(Z_{i,t})\beta(Z_{i,t})\right\|^2
\right]
+\frac{\sigma_\varepsilon^2}{N}\mathbb E_\varepsilon[f(Z)^2].
\end{aligned}
\]

Define the kernel covariance operator
\[
T_K=\mathbb E_\varepsilon[K(\cdot,Z)\otimes K(\cdot,Z)],
\]
where $Z\sim P_Z$. By the reproducing property,
\[
\langle f,T_Kf\rangle_{\mathcal H_K}
=\mathbb E_\varepsilon[f(Z)^2].
\]
Since the first term in the preceding expression for $\langle f,\Sigma_{\mathcal H_K,\varepsilon}f\rangle_{\mathcal H_K}$ is nonnegative, we immediately obtain
\[
\langle f,\Sigma_{\mathcal H_K,\varepsilon}f\rangle_{\mathcal H_K}
\geq\frac{\sigma_\varepsilon^2}{N}
\langle f,T_Kf\rangle_{\mathcal H_K}.
\]

To establish the upper bound, let $B_t\in\mathbb R^{N\times3}$ denote the matrix whose $i$-th row is $\beta(Z_{i,t})^\top$, and write $f_t=(f(Z_{1,t}),\ldots,f(Z_{N,t}))^\top$. The cross-sectional balance condition implies
\[
B_t^\top B_t=NgI_3.
\]
Consequently, the largest eigenvalue of $B_tB_t^\top$ is $Ng$, and hence
\[
\begin{aligned}
\left\|\frac1N\sum_{i=1}^N f(Z_{i,t})\beta(Z_{i,t})\right\|^2
&=\frac1{N^2}f_t^\top B_tB_t^\top f_t\\
&\leq\frac{g}{N}\|f_t\|^2
=\frac{g}{N}\sum_{i=1}^N f(Z_{i,t})^2.
\end{aligned}
\]
Taking expectations and using the common marginal law,
\[
\mathbb E_\varepsilon\left[
\left\|\frac1N\sum_{i=1}^N f(Z_{i,t})\beta(Z_{i,t})\right\|^2
\right]\leq g\mathbb E_\varepsilon[f(Z)^2].
\]
Combining this inequality with the idiosyncratic component yields
\[
\langle f,\Sigma_{\mathcal H_K,\varepsilon}f\rangle_{\mathcal H_K}
\leq\left(g+\frac{\sigma_\varepsilon^2}{N}\right)
\langle f,T_Kf\rangle_{\mathcal H_K}.
\]
Therefore, for every $f\in\mathcal H_K$,
\[
\frac{\sigma_\varepsilon^2}{N}
\langle f,T_Kf\rangle_{\mathcal H_K}
\leq
\langle f,\Sigma_{\mathcal H_K,\varepsilon}f\rangle_{\mathcal H_K}
\leq
\left(g+\frac{\sigma_\varepsilon^2}{N}\right)
\langle f,T_Kf\rangle_{\mathcal H_K}.
\]
Equivalently, in the Loewner ordering of self-adjoint operators,
\[
\frac{\sigma_\varepsilon^2}{N}T_K
\preceq\Sigma_{\mathcal H_K,\varepsilon}
\preceq
\left(g+\frac{\sigma_\varepsilon^2}{N}\right)T_K.
\]

Let $\lambda_j(T_K)$ denote the eigenvalues of $T_K$, arranged in nonincreasing order. Since both operators are positive, self-adjoint and compact, the min--max principle implies
\[
\frac{\sigma_\varepsilon^2}{N}\lambda_j(T_K)
\leq\mu_{j,\varepsilon}
\leq
\left(g+\frac{\sigma_\varepsilon^2}{N}\right)\lambda_j(T_K),
\qquad j\geq1.
\]

For a Matérn kernel with smoothness parameter $\nu$ on the bounded Lipschitz domain $\mathcal Z=[-1,1]^D$, equipped with a uniform probability measure, the standard spectral asymptotics of the associated kernel integral operator imply that there exist constants $0<c_K\leq C_K<\infty$ such that
\[
c_Kj^{-(1+2\nu/D)}
\leq\lambda_j(T_K)
\leq C_Kj^{-(1+2\nu/D)},
\qquad j\geq1.
\]
The nonzero eigenvalues of the kernel integral operator on $L^2(P_Z)$ coincide with those of $T_K$ on $\mathcal H_K$. Therefore, combining the preceding inequalities, we obtain
\[
\frac{\sigma_\varepsilon^2}{N}c_Kj^{-(1+2\nu/D)}
\leq\mu_{j,\varepsilon}
\leq
\left(g+\frac{\sigma_\varepsilon^2}{N}\right)
C_Kj^{-(1+2\nu/D)}.
\]
Thus the required polynomial eigenvalue decay holds with
\[
b=1+\frac{2\nu}{D},
\qquad
c_\mu=\frac{\sigma_\varepsilon^2}{N}c_K,
\qquad
C_\mu=\left(g+\frac{\sigma_\varepsilon^2}{N}\right)C_K.
\]
In the baseline specification, $\nu=3/2$ and $D=6$, so that
\[
b=1+\frac{2(3/2)}6=\frac32.
\]
Hence,
\[
c_\mu j^{-3/2}\leq\mu_{j,\varepsilon}\leq C_\mu j^{-3/2},
\qquad j\geq1,
\]
which establishes the claimed eigenvalue decay.
\end{proof}

\begin{proposition}[Uniform bounds for the optimal Sharpe ratio]
Consider the factor economy described above, with $N=3G\geq3$, $g=c_\beta/3>0$, $\sigma_\varepsilon^2>0$, and $0<\|\mu_F\|^2<1$. Let
\[
W^\star(z)=\frac{\beta(z)^\top\mu_F}{g+\sigma_\varepsilon^2/N}
\]
denote the population-optimal portfolio policy. Then there exist constants $0<SR_{\min}\leq SR_{\max}<\infty$, independent of $N$, such that
\[
SR_{\min}\leq SR(W^\star)\leq SR_{\max}.
\]
\end{proposition}

\begin{proof}
Recall that the portfolio return associated with a policy $W$ is
\[
R^p_{t+1}(W)=\frac1N\sum_{i=1}^N W(Z_{i,t})R_{i,t+1}.
\]
Its Sharpe ratio is defined as
\[
SR(W)=\frac{\mathbb E[R^p_{t+1}(W)]}
{\sqrt{\operatorname{Var}(R^p_{t+1}(W))}}.
\]

For the population-optimal policy, the cross-sectional balance condition implies
\[
b_t(W^\star)=\frac1N\sum_{i=1}^N W^\star(Z_{i,t})\beta(Z_{i,t})
=\frac{g}{g+\sigma_\varepsilon^2/N}\mu_F.
\]
Consequently,
\[
\mathbb E[R^p_{t+1}(W^\star)]
=\mu_F^\top b_t(W^\star)
=\frac{g\|\mu_F\|^2}{g+\sigma_\varepsilon^2/N}.
\]

Using the conditional second-moment assumptions of the factor model, we also have
\[
\begin{aligned}
\mathbb E[(R^p_{t+1}(W^\star))^2]
&=\mathbb E[\|b_t(W^\star)\|^2]
+\frac{\sigma_\varepsilon^2}{N}\mathbb E[W^\star(Z)^2]\\
&=\frac{g^2\|\mu_F\|^2}{(g+\sigma_\varepsilon^2/N)^2}
+\frac{\sigma_\varepsilon^2}{N}
\frac{g\|\mu_F\|^2}{(g+\sigma_\varepsilon^2/N)^2}\\
&=\frac{g\|\mu_F\|^2}{g+\sigma_\varepsilon^2/N},
\end{aligned}
\]
where we used $\mathbb E[\beta(Z)\beta(Z)^\top]=gI_3$.

Define
\[
q_N=\frac{g\|\mu_F\|^2}{g+\sigma_\varepsilon^2/N}.
\]
The preceding identities imply
\[
\mathbb E[R^p_{t+1}(W^\star)]
=\mathbb E[(R^p_{t+1}(W^\star))^2]=q_N,
\]
and therefore
\[
\operatorname{Var}(R^p_{t+1}(W^\star))=q_N-q_N^2.
\]
It follows that
\[
SR(W^\star)=\frac{q_N}{\sqrt{q_N-q_N^2}}
=\sqrt{\frac{q_N}{1-q_N}}.
\]

Since $N\geq3$, we have
\[
0<\frac{\sigma_\varepsilon^2}{N}\leq\frac{\sigma_\varepsilon^2}{3}.
\]
Together with $g>0$ and $0<\|\mu_F\|^2<1$, this yields
\[
0<
\frac{g\|\mu_F\|^2}{g+\sigma_\varepsilon^2/3}
\leq q_N
\leq\|\mu_F\|^2<1.
\]
Since the function $q\mapsto\sqrt{q/(1-q)}$ is strictly increasing on $(0,1)$, we obtain
\[
\sqrt{\frac{g\|\mu_F\|^2}
{g+\sigma_\varepsilon^2/3-g\|\mu_F\|^2}}
\leq SR(W^\star)
\leq
\sqrt{\frac{\|\mu_F\|^2}{1-\|\mu_F\|^2}}.
\]

Thus, defining
\[
SR_{\min}=
\sqrt{\frac{g\|\mu_F\|^2}
{g+\sigma_\varepsilon^2/3-g\|\mu_F\|^2}},
\qquad
SR_{\max}=
\sqrt{\frac{\|\mu_F\|^2}{1-\|\mu_F\|^2}},
\]
we conclude that
\[
0<SR_{\min}\leq SR(W^\star)\leq SR_{\max}<\infty.
\]
Both constants depend only on the fixed model parameters $g$, $\mu_F$, and $\sigma_\varepsilon^2$, and are independent of the number of assets $N$. This establishes the required uniform Sharpe-ratio bounds.
\end{proof}
